
\definecolor{myBLUE}{RGB}{72,88,215}
\definecolor{myYELLOW}{RGB}{249,227,71}
\definecolor{myRED}{RGB}{201,59,84}
\definecolor{myPERPULE}{RGB}{84,29,149}
\definecolor{myGREEN}{RGB}{90,142,99}
\definecolor{myORANGE}{RGB}{222,118,3}

\definecolor{myYELLOWGREEN}{RGB}{160,178,71}
\definecolor{myREDORANGE}{RGB}{205,66,0}
\definecolor{myPERPULERED}{RGB}{117,46,102}
\definecolor{myORANGEYELLOW}{RGB}{235,171,27}
\definecolor{myGREENBLUE}{RGB}{60,99,104}
\definecolor{myBLUEPURPULE}{RGB}{88,75,176}





\definecolor{CNblack}{RGB}{0,0,0}
\definecolor{CNblue}{RGB}{0,0,255}
\definecolor{CNbrown}{RGB}{127.5,102,63.75}
\definecolor{CNgrey}{RGB}{127.5,127.5,127.5}
\definecolor{CNgreen}{RGB}{0,255,0}
\definecolor{CNorange}{RGB}{255,204,0}
\definecolor{CNpink}{RGB}{255,127.5,255}
\definecolor{CNpurple}{RGB}{255,0,255}
\definecolor{CNred}{RGB}{255,0,0}
\definecolor{CNwhite}{RGB}{255,255,255}
\definecolor{CNyellow}{RGB}{255,255,0}







% latin abbreviations
\newcommand{\latin}[1]{\textit{#1}}
\newcommand\eg{\textit{e.g.}\xspace}
\newcommand\ie{\textit{i.e.}\xspace}
\newcommand\etal{\latin{et al.}\xspace}

% names of sections, equations, ...
\newcommand\equationname{Eq.}
\newcommand\sectionname{Sect.}
\newcommand\chapternname{Chap.}

\newcommand\wrt{with respect to\@\xspace}


% references to ...
% ...figures
\newcommand{\figref}[1]{\figurename~\ref{#1}}%{Fig.~\ref{#1}}

% ...tables
\newcommand{\tabref}[1]{\tablename~\ref{#1}}%{Tab.~\ref{#1}}

% ...equations
\newcommand{\eqnref}[1]{\equationname~\eqref{#1}}%{Eq.~\eqref{#1}}

% ...chapters
\newcommand{\chapref}[1]{\chaptername~\ref{#1}}%{Chap.~\ref{#1}}

% ...sections
\newcommand{\secref}[1]{\sectionname~\ref{#1}}%{Sec.~\ref{#1}}

% ...appendices
\newcommand{\appref}[1]{Appendix~\ref{#1}}%{App.~\ref{#1}}

% ...definitions
\newcommand{\defref}[1]{Definition~\ref{#1}}%{Def.~\ref{#1}}

% ...algorithms
\newcommand{\algoref}[1]{Algorithm~\ref{#1}}%{Alg.~\ref{#1}}

% ... theorems in appendix
\newtheorem{theorem}{Theorem}

% to ``highlight'' some parts of an equation
\newcommand\partOfExpression[1]{\langle\mathfrak{#1}\rangle}



%  we always want to place a nice intro question or quote at the start of every chapter...
\newcommand\chapStartSentence[1]{ \begin{center}\textit{#1}\end{center}}


%your own macros, abbreviations, etc.

%\newcommand\todo[1]{\textcolor{red}{TODO: #1}}
\newcommand\fixme[1]{\textcolor{green}{FIXME: #1}}


\newcommand{\keyword}[1]{\emph{#1}}

\newcommand{\ROT}[1]{\keyword{Rule of thirds}\index{Rule of thirds} }

